% Chapter 3. Source pages PDF 73--90 (printed 65--82).
\section{Constructions with sets}
\begin{definition}[Cardinality comparison]
\label{def:lbc-cardinality}
\source{leinster-basic-category-theory:page-0082}
For sets $A,B$, write $|A|\le |B|$ when an injection $A\to B$ exists,
$|A|\ge|B|$ dually, and $|A|=|B|$ when a bijection exists.
\end{definition}

\begin{definition}[Initial and terminal sets]
\label{def:lbc-empty-singleton}
\source{leinster-basic-category-theory:page-0075}
The empty set is initial in $\Set$, and a singleton is terminal.
\uses{def:lbc-initial-terminal}
\end{definition}

\begin{remark}[Sets form a category]
\label{rem:lbc-set-category}
\source{leinster-basic-category-theory:page-0075}
Sets and functions form the category $\Set$, with ordinary composition,
associativity, and identity functions.
\uses{def:lbc-category}
\end{remark}

\begin{definition}[Products and sums of sets]
\label{def:lbc-set-products-sums}
\source{leinster-basic-category-theory:page-0076}
Any sets $A,B$ have a product $A\times B$ of ordered pairs and a disjoint sum
$A+B$. More generally, a family $(A_i)_{i\in I}$ has product
$\prod_iA_i$ of choice families and disjoint sum $\coprod_iA_i$.
\uses{def:lbc-category}
\end{definition}

\begin{definition}[Sets of functions]
\label{def:lbc-set-function-set}
\source{leinster-basic-category-theory:page-0077}
For sets $A,B$, the set $A^B$ consists of functions $B\to A$ and is the
product of the constant family $(A)_{b\in B}$. For finite sets, products, sums,
and function sets have cardinalities $mn$, $m+n$, and $m^n$; the usual
distributive and exponential isomorphisms hold for all sets.
\uses{def:lbc-set-products-sums}
\end{definition}

\begin{definition}[Equalizer and coequalizer in sets]
\label{def:lbc-set-equalizer}
\source{leinster-basic-category-theory:page-0078}
For functions $f,g:A\to B$, the equalizer is
$\{a\in A\mid f(a)=g(a)\}\hookrightarrow A$; the coequalizer is the quotient
of $B$ by the least equivalence relation identifying $f(a)$ with $g(a)$.
\end{definition}

\begin{definition}[Power set and cardinal arithmetic]
\label{def:lbc-power-set}
\source{leinster-basic-category-theory:page-0077}
The power set $\mathcal P(A)$ is the set of all subsets of $A$; cardinal sums
and products are represented by disjoint union and cartesian product.
\end{definition}

\begin{remark}[Subsets and characteristic functions]
\label{rem:lbc-characteristic-functions}
\source{leinster-basic-category-theory:page-0077}
For the two-element set $2=1+1$, subsets $S\subseteq A$ correspond bijectively
to characteristic functions $\chi_S:A\to2$; the power set $\mathcal P(A)$ is
the set $2^A$ of such functions.
\uses{def:lbc-power-set,def:lbc-set-function-set}
\end{remark}

\begin{remark}[Quotients of sets]
\label{rem:lbc-set-quotient}
\source{leinster-basic-category-theory:page-0078}
Given an equivalence relation $\sim$ on $A$, the quotient map
$p:A\to A/{\sim}$ is surjective and satisfies
$p(a)=p(a')\iff a\sim a'$. Maps $A/{\sim}\to B$ correspond uniquely to maps
$A\to B$ constant on equivalence classes.
\uses{def:lbc-set-equalizer}
\end{remark}

\begin{remark}[Natural numbers and choice]
\label{rem:lbc-natural-numbers-choice}
\source{leinster-basic-category-theory:page-0079}
The natural numbers carry $0$ and successor $s$ with the recursion property:
for every set $X$, $a\in X$, and $r:X\to X$, there is a unique
$x:\mathbb N\to X$ with $x(0)=a$ and $x\circ s=r\circ x$. A section of
$f:A\to B$ is a right inverse; the axiom of choice asserts that every
surjection of sets has a section.
\uses{def:lbc-category}
\end{remark}

\section{Small and large categories}
\begin{definition}[Locally small and essentially small]
\label{def:lbc-locally-small}
\source{leinster-basic-category-theory:page-0083}
A category is locally small if every hom-class is a set. It is essentially small
if it is equivalent to a small category.
\uses{def:lbc-category,def:lbc-equivalence}
\end{definition}

\begin{theorem}[Cantor--Bernstein]
\label{thm:lbc-cantor-bernstein}
\source{leinster-basic-category-theory:page-0082}
\dcref{3.2.1}
If $|A|\le|B|\le|A|$, then $A\cong B$.
\notready
\end{theorem}
\begin{proof}
The source defers this proof to Exercise 3.2.12.
\end{proof}

\begin{theorem}[Cantor]
\label{thm:lbc-cantor}
\source{leinster-basic-category-theory:page-0082}
\dcref{3.2.2}
For every set $A$, $|A|<|\mathcal P(A)|$.
\uses{def:lbc-power-set}
\end{theorem}
\begin{proof}
The source defers this proof to Exercise 3.2.13.
\end{proof}

\begin{corollary}[Unbounded cardinalities]
\label{cor:lbc-unbounded-cardinality}
\source{leinster-basic-category-theory:page-0082}
\dcref{3.2.3}
For every set $A$ there is a set $B$ with $|A|<|B|$.
\uses{thm:lbc-cantor}
\end{corollary}
\begin{proof}
Take $B=\mathcal P(A)$ and apply Theorem~\ref{thm:lbc-cantor}.
\end{proof}

\begin{proposition}[A set outside a given family]
\label{prop:lbc-set-outside-family}
\source{leinster-basic-category-theory:page-0083}
\dcref{3.2.4}
For a set $I$ and a family of sets $(A_i)_{i\in I}$, there exists a set not
isomorphic to any of the $A_i$.
\uses{thm:lbc-cantor}
\end{proposition}
\begin{proof}
Put $A=\mathcal P(\coprod_{i\in I}A_i)$. For each $j$, the inclusion
$A_j\hookrightarrow\coprod_iA_i$ and Cantor give $|A_j|<|A|$, so
$A_j\not\cong A$.
\end{proof}

\begin{proposition}[Set is not essentially small]
\label{prop:lbc-set-not-essentially-small}
\source{leinster-basic-category-theory:page-0084}
\dcref{3.2.8}
The category $\Set$ is not essentially small.
\uses{thm:lbc-cantor}
\end{proposition}
\begin{proof}
If $\Set$ were equivalent to a small category, choose one representative of
each isomorphism class. Their underlying sets form a set $U$ containing a set
of every cardinality, contradicting Cantor applied to $\mathcal P(U)$.
\end{proof}

\begin{definition}[Category of small categories]
\label{def:lbc-cat}
\source{leinster-basic-category-theory:page-0085}
\dcref{3.2.10}
The category $\Cat$ has small categories as objects and functors as maps.
\uses{def:lbc-category,def:lbc-functor}
\end{definition}
